% Proposed repair of the original H5 arithmetic step. Independent review pending.
% Uses the existing baseline bibliography keys, plus GaudronRemond2023 and
% the published 2021 edition of Voight, chapter 43.
\subsection{A level bound for the non-CM fibers}
\label{sec:h5-level-repair}

This argument allows the scalar character order and coefficient field to
vary. It does not assert bounded same-field descent to a rational
GL$_2$-type variety. Let $s$ be a rational point outside the five contacts,
let $L/\mathbb Q$ be the normal fiber field constructed above, and let
$A/L$ be a non-CM fiber surface. Thus $[L:\mathbb Q]\mid576$, all
geometric homomorphisms between conjugate fibers are defined over $L$,
and $\operatorname{End}^0(A_{\overline{\mathbb Q}})=B$, the quaternion
division algebra of discriminant $210$. Put
\[
 R=B_*M(s),\qquad T=\{p:p\mid R\}.
\]
The geometric argument gives that $L$ is unramified outside $T$ and
$A$ has good reduction outside the primes above $210$. All conjugate
QM structures lie in the same normalizer orbit.

\paragraph{Good reduction after full torsion.}
First a QM surface with the indicated division action has potentially
good reduction at every finite place. Indeed, pass to an extension with
semistable reduction, using \cite[Theorem 3.4, p.~406]{silverbergzarhin}.
Every endomorphism extends uniquely to the N\'eron model by its mapping
property. On the identity component of the geometric special fiber it
preserves the maximal torus $\mathcal T$: this torus is the maximal
connected affine subgroup, and every homomorphism from a torus to an
abelian variety is zero. If its rank $t$ is positive, restriction gives
a unital homomorphism
\[
 B\longrightarrow\operatorname{End}^0(\mathcal T)=M_t(\mathbb Q),
 \qquad 1\le t\le2.
\]
The kernel is a two-sided ideal of the simple algebra $B$, and the map
is unital, so it is injective. For $t=1$ dimensions prohibit this;
for $t=2$ it would be an isomorphism from a division algebra to
$M_2(\mathbb Q)$, also impossible. Hence $t=0$. The semistable identity
component is abelian, which is good reduction. These N\'eron-model and
reduction conventions are those of \cite[\S3, pp.~405--406]{silverbergzarhin}.

Set $E=L(A[11])$. We next justify that $E/\mathbb Q$ is Galois.
For each $\sigma$, the normalizer relation supplies a compatible
isogeny ${}^{\sigma}A\to A$ of degree prime to $11$. Here is the
prime-to-$11$ detail. In the analytic description
$A_z=\mathbb C^2/(\mathcal O(z,1)^t)$ of the published 2021 edition
of Voight \cite[equations (43.1.2)--(43.1.3)]{voight-qm2021}, a positive normalizer element
$b=\left(\begin{smallmatrix}a&b'\\c&d\end{smallmatrix}\right)$
gives the scalar compatible map $(cz+d)I:A_{bz}\to A_z$.
When $\mathcal O b\subset\mathcal O$, its degree is
$[\mathcal O:\mathcal O b]=\operatorname{Nrd}(b)^2$.
At $11$ the algebra is split and the maximal-order normalizer is
$\mathbb Q_{11}^*\operatorname{GL}_2(\mathbb Z_{11})$.
Rescale $b$ rationally to make it an $11$-adic unit, then clear the
remaining denominators by an integer prime to $11$. This produces
the claimed isogeny. Changes of level marking use isomorphisms.
All these geometric homomorphisms are $L$-defined by construction.
Consequently the $11$-torsion representations of $A$ and its conjugates
over $L$ have the same kernel. Thus every conjugate of $E$ equals $E$.

The faithful action on $A[11]$ now gives
\[
 [E:\mathbb Q]\le576|\operatorname{GL}_4(\mathbb F_{11})|
 <2^{10}11^{16}<2^{66}.
\tag{LR1}
\]
For residue characteristic different from $11$, full rational
$11$-torsion gives semistable reduction over $E$ by
\cite[Theorem 3.5, p.~406]{silverbergzarhin}. The preceding torus
argument makes this good reduction. In particular it applies at
each of $2,3,5,7$.

\paragraph{The scalar lift with controlled support.}
Choose $B$-compatible isogenies $\mu_\sigma:{}^\sigma A\to A$ over $L$,
with $\mu_1=1$, and inflate this finite choice to $G_{\mathbb Q}$.
Their discrepancies $c(\sigma,\tau)$ commute with $B$, so lie in
$\mathbb Q^*$. Put $a_\sigma=(\deg\mu_\sigma)^{-1/4}>0$ and
$\nu_\sigma=a_\sigma\mu_\sigma$, with trivial Galois action on the
added algebraic coefficients. Since a scalar $c$ has degree $c^4$,
the discrepancies of the $\nu_\sigma$ form a sign cocycle
$\varepsilon\in Z^2(G_{\mathbb Q},\{\pm1\})$, inflated from
$\operatorname{Gal}(L/\mathbb Q)$.

The sign class is unramified, hence zero, locally outside $T$.
At each odd prime where its Brauer invariant is $1/2$, choose an odd
Dirichlet character modulo $p$; at $2$ use the odd character modulo
$4$ if required. Their product $\chi$ is tame at every odd prime,
has local conductor at most $4$ at $2$, and has order
\[
 m\le2\prod_{\substack{p\in T\\p\ne2}}(p-1)\le2R.
\]
For constant coefficient groups, the boundary in
$1\to\{\pm1\}\to\mu_\infty\xrightarrow{2}\mu_\infty\to1$
is the obstruction to a character square root. Local reciprocity
identifies this obstruction with the value at $-1$: the cyclic
torsion factor has precisely this obstruction, while the valuation
and principal-unit factors admit character square roots.
Thus the boundary of $\chi$ equals $[\varepsilon]$ at every finite
place. Brauer reciprocity and the global-to-local injection give
equality globally. A pointwise square root, corrected by a sign
cochain, gives
\[
 b:G_{\mathbb Q}\longrightarrow\mu_{2m},\qquad
 b^2=\chi,\qquad\delta b=\varepsilon.
\tag{LR2}
\]
Only this \emph{sign} cochain is asserted to take root-of-unity values;
the full scalar splitting also involves the $a_\sigma$.

On $G_L$, $b$ is a conjugation-invariant character. Outside $T$ its
inertia characters are quadratic. Only finitely many are nontrivial;
twist by the quadratic character of the product of those odd primes
to kill them. At each odd prime there is a unique nontrivial quadratic
inertia character. This changes no square or coboundary in (LR2),
and introduces possible ramification only at $2\in T$.
The kernel field $L'$ of $b|G_L$ is normal over $\mathbb Q$ and
\[
 [L':\mathbb Q]\le2m[L:\mathbb Q]\le2304R.
\tag{LR3}
\]

\paragraph{The compatible modular representations.}
Choose a number field containing all normalization scalars, the values
of $b$, and a splitting field of $B$. For each coefficient prime $\ell$
and an embedding into $\overline{\mathbb Q}_\ell$, set
\[
 \rho_\ell(\sigma)=b(\sigma)^{-1}a_\sigma\mu_\sigma\circ\sigma.
\]
The cocycle identity makes this a representation commuting with $B$;
after splitting $B$, matrix units give $\rho_\ell=r_\ell\oplus r_\ell$
with $\dim r_\ell=2$. On $G_{L'}$, $\rho_\ell$ is the Tate action of
$A$. Faltings' semisimplicity and endomorphism theorems and $\operatorname{End}^0(A)=B$
show that the multiplicity representation is absolutely irreducible.
Proper-smooth comparison and finite descent, as in
Lemma~\ref{lem:de-rham-descent}, make $r_\ell$ de Rham with Tate
Hodge--Tate weights $\{0,-1\}$; Berger's monodromy theorem makes it
potentially semistable.

It is odd. On Betti homology complex conjugation followed by the
holomorphic isogeny $\mu_c$ exchanges the two Hodge summands, each
a simple $B\otimes\mathbb C$-module. Its action on the multiplicity
space is therefore nonscalar. After the scalar correction its square
is one, so its eigenvalues are $1,-1$. The representation is unramified
outside $T\cup\{\ell\}$, by (LR2) and the reduction and field properties.
Pan's published Theorem 1.0.4 applies at $\ell=5$, with no additional
residual-image restriction at that prime, and identifies $r_5$ with
a cuspidal modular representation up to a character twist.

The twist can be taken finite and the weight is two, without invoking
a separate local-algebraicity theorem for characters. The twisting
character is a de Rham subrepresentation of the relevant Hom
representation. Comparing its integer weight with $\{0,-1\}$ and
the modular weights $\{0,1-k\}$ forces $k=2$ and twist weight zero.
On $G_{L'}$, the identity
$(\det r_5)^2=\det V_5(A)=\chi_{\rm cyc}^2$ shows that
$\det r_5/\chi_{\rm cyc}$ has finite image, since $G_{L'}$ has finite
index in $G_{\mathbb Q}$. The modular determinant
divided by $\chi_{\rm cyc}$ also has finite image. The square of the
twisting character consequently has finite image, so the character
itself has finite image. Absorb it in the modular form and take the
primitive weight-two newform $f$ of level $N$.

The conductor at $5$ is not omitted. For each prime
$p\notin T\cup\{3,5\}$, reduction of $\mu_\sigma$ composed with
relative $p$-Frobenius is a quasi-endomorphism of the good reduction
of $A$ at a chosen place. After clearing denominators its characteristic
polynomial is independent of the Tate prime. Half its trace, multiplied
by $b(\sigma)^{-1}a_\sigma$, is therefore one algebraic number giving
the trace of every $r_\ell$. Enlarge the coefficient field to contain
the coefficients of $f$. The equality with the trace of $f$ at five
is an equality of algebraic numbers; transporting it to an embedding
at three and applying Chebotarev and semisimplicity identifies $r_3$
with the three-adic realization of $f$. Thus use $r_3$ to compute the
conductor at $5$ and $r_5$ at every other prime. The modular level is
the conductor of these realizations by
\cite[\S\S0.3--0.7, Theorem A, pp.~409--411]{carayol}
and newform theory \cite[Theorems 1 and 4]{casselman}.
Here the totally real field is $\mathbb Q$, of odd degree one, so
Carayol's extra finite-place hypothesis for even degree does not apply.
His unramified normalization twist and contragredient preserve the
conductor; his theorem allows arbitrary nebentypus.
Li's Fricke relation and Shimura's
Jacobian construction, as in the baseline modularity argument, also
place $r_5$ in the Tate module of $J_1(N)$.

\paragraph{The conductor bound.}
Let $F=E\mathbb Q(\chi)$, where $\mathbb Q(\chi)$ means the kernel
field of the character, not its value field. On $G_F$, $b$ is a
conjugation-invariant quadratic character. Its kernel gives a Galois
extension $F'/\mathbb Q$ of degree at most two over $F$.
For $p=2,3,5,7$, let $e_E,e_{F'}$ denote local ramification indices
over $\mathbb Q_p$. From (LR1), $v_p(e_E)\le65$.
At odd $p$ both added extensions are tame, whereas at $2$ each
adds ramification index at most two. Consequently
\[
 v_p(e_{F'})\le65\ (p=3,5,7),\qquad v_2(e_{F'})\le67.
\tag{LR4}
\]
The possibly unbounded tame ramification index causes no problem.
For a finite Galois extension of $\mathbb Q_p$ with ramification
index $e$ and different exponent $\delta$, an Eisenstein polynomial
over the maximal unramified subfield gives
\[
 \delta\le e-1+ev_p(e).
\]
Indeed the derivative terms have distinct valuations modulo $e$,
and its leading derivative term has the stated valuation. In upper
numbering the different formula is
\[
 \delta/e=1-1/e+\int_0^\infty(1-|G^u|^{-1})\,du.
\]
Before the last positive break the integrand is at least $1/2$.
Thus (LR4) bounds this break for $F'/\mathbb Q_p$ by $136$.
These ramification and conductor conventions are those of
\cite[\S\S4.1--4.3, pp.~155--159]{serrelocal}.
Over $F'$ the scalar correction is trivial and $A$ has good reduction.
Above that break, therefore, the two-dimensional representation is
trivial on upper wild groups. Its Swan conductor is at most $272$,
and its tame conductor is at most two. Hence
\[
 a_p(r)\le274<300\qquad(p=2,3,5,7).
\]
For $p>7$, $L/\mathbb Q$ is tame because its degree divides $576$,
$A/L$ has good reduction, and $b$ is trivial on wild inertia:
on that pro-$p$ group its square is tame and its remaining image has
order at most two. Thus $a_p(r)\le2$ there; it is zero outside $T$.
Using the coefficient-prime comparison above proves
\[
 N\le210^{300}R^2.
\tag{LR5}
\]

\paragraph{The original constant is preserved.}
Apply the accompanying modular-Jacobian transfer calculation with
$C=210^{300}$ and (LR3)--(LR5). It gives
\[
 h_F(A)<2^{24}210^{3600}R^{24}
 <2^{28824}R^{24}<4608^{165888}R^{24}.
\tag{LR6}
\]
The last comparison uses $210<2^8$, $4608>2^{12}$, and
$28824<1990656=12\cdot165888$. Together with the separately proved
CM bound and geometric height comparison, (LR6) yields exactly
\[
 h(s)\le
 \bigl(2(4608^{165888}+6^{100})B_*^{24}+128\bigr)M(s)^{24}.
\]
No contact, coefficient constant, or radical exponent has been changed.
The proposed repair must be independently reviewed before replacing
the conditional status of the frozen H5 snapshot.
